% NexusQuant — Abstract
% NeurIPS 2026 submission
% All numbers GPU-validated on A10G. Ranges reported honestly.

\begin{abstract}
The key-value (KV) cache is the primary memory bottleneck during autoregressive inference of large language models, accumulating over 40\,GB at 128K tokens for a 70B-parameter model. Existing training-free methods achieve either high compression with significant quality loss through quantization, or moderate compression with near-zero loss through token eviction, but rarely combine both effectively. We present \textsc{NexusQuant}, a training-free, zero-calibration pipeline that combines attention-aware token eviction with E8 lattice vector quantization, Hadamard incoherence preprocessing, temporal delta coding, and zstd entropy compression. All compression ratios are measured with full overhead included. On Mistral-7B evaluated on multi-topic 4K-token passages, \textsc{NexusQuant} achieves 10.4x compression at $+0.14\%$ perplexity degradation with 35\% eviction, 16.6x at $+0.82\%$ with 60\% eviction, and up to 32.7x at $+2.1\%$ for aggressive configurations. A central empirical finding is context-length scaling, whereby quality improves 5--6x as prefix length grows from 500 to 2048 tokens, because longer contexts provide more signal for the attention-based importance scorer to distinguish important tokens from noise. On short-context inputs of approximately 500 tokens, the safe operating point is 10x at less than 1\% PPL degradation. Results are consistent across academic, technical, and creative text domains, with creative narrative being the harder case at low eviction rates. We additionally validate two algorithmic improvements on A100-SXM4-40GB hardware. First, replacing the key-key proxy scorer with real accumulated softmax attention weights during prefill (the \emph{real attention scorer}) reduces quality loss at 80\% eviction from $+3.20\%$ to $+0.66\%$ PPL---a 4.8x improvement---so that 80\% eviction with the real scorer matches 35\% eviction with the proxy scorer; this advance requires \texttt{eager} attention mode and one extra prefill pass. Second, asymmetric quantization using 3-bit keys and 2-bit values (K3V2) at 35\% eviction yields $+0.35\%$ PPL at $8.89\times$ compression, compared to $+2.26\%$ PPL at $10.52\times$ for symmetric 2-bit (K2V2)---a 6.4x quality improvement at 15\% ratio cost; the asymmetry is justified because key quantization noise shifts all softmax attention weights, whereas value noise only scales outputs proportionally. Combining K3V2 with the real scorer targets sub-$9\times$ compression at $<$0.5\% PPL. Compared to published methods, \textsc{NexusQuant} exceeds TurboQuant by 1.7--2.8x compression at comparable quality, surpasses CommVQ and Palu on both ratio and quality, and approaches KVTC without requiring any calibration data. Main results are on Mistral-7B. Llama-3-8B cross-architecture results exist but exhibit an evaluation artifact from attention mask renormalization. Evaluation uses perplexity on multi-topic passages and five QA tasks; full LongBench task accuracy is not yet reported.
\end{abstract}
